% GENERATED FILE. Edit canonical lessons, then run npm run generate:books.
% canonical-source-hash: fnv1a64:3303f7783c104ad8
% canonical-lessons: PA-C321-zukam, PA-C321-khangh, PA-C321-chhikk, PA-C321-sah, PA-C321-nigalna

\chapter{A Cough and a Cold}
\label{ch:pa-a2-a-cough-and-a-cold}
\hlchaptermodality{321}

\begin{chapteropening}
\textbf{By the end of this chapter:} \emph{I can say that I have a cold or a cough.}

Complete the last lesson of chapter 321: I can say that I have a cold or a cough.
\end{chapteropening}

\section[\texorpdfstring{\pa{ਜ਼ੁਕਾਮ} (zukām)}{ਜ਼ੁਕਾਮ (zukām)}]{\pa{ਜ਼ੁਕਾਮ} (zukām) — a head cold}

\label{lesson:PA-C321-zukam}



\begin{quote}
Before the new one: say the Punjabi for the pulse, then the Punjabi for an injection.
\end{quote}

\subsection*{You'll want to know: \pa{ਜ਼ੁਕਾਮ}}
\textbf{\pa{ਜ਼ੁਕਾਮ}} — \emph{zukām} — \textquotedblleft{}a head cold\textquotedblright{}.

\begin{quote}
\pa{ਮੈਨੂੰ} \pa{ਜ਼ੁਕਾਮ} \pa{ਹੈ} \pa{ਅਤੇ} \pa{ਮੈਂ} \pa{ਅੱਜ} \pa{ਘਰ} \pa{ਵਿੱਚ} \pa{ਹਾਂ।}
\end{quote}

\begin{quote}
\emph{mainū̃ zukām hai ate maiṁ ajj ghar vich hāṁ.}
\end{quote}

\begin{quote}
I have a cold, and I am at home today.
\end{quote}

A masculine noun.

\begin{cousinweb}
\textbf{\pa{ਜ਼ੁਕਾਮ}} (\emph{zukām}) is a loan, usually traced through Persian to Arabic \emph{zukām}, \textquotedblleft{}a cold, catarrh\textquotedblright{}.
\end{cousinweb}

\begin{grammarlens}[title={What you've built}]
One more word for this chapter.
\end{grammarlens}

\subsection*{Your turn}
\begin{itemize}
\raggedright
  \item \emph{Say it:} \emph{zukām}
  \item \emph{Say it:} \emph{zukām}, once more
  \item \emph{Say it:} \emph{mainū̃ zukām hai ate maiṁ ajj ghar vich hāṁ.} — I have a cold, and I am at home today.
  \item \emph{Recall:} say the Punjabi for the pulse, then the Punjabi for an injection, then say \emph{zukām} again
\end{itemize}

\begin{tcolorbox}[breakable,colback=teal!4,colframe=teal!35!black,title={Before you move on}]
What does \pa{ਜ਼ੁਕਾਮ} mean? (\textquotedblleft{}A head cold\textquotedblright{}.) Say it once more.
\end{tcolorbox}
\section[\texorpdfstring{\pa{ਖੰਘ} (khaṅgh)}{ਖੰਘ (khaṅgh)}]{\pa{ਖੰਘ} (khaṅgh) — a cough}

\label{lesson:PA-C321-khangh}



\begin{quote}
Before the new one: say the Punjabi for an injection, then the Punjabi for a head cold.
\end{quote}

\subsection*{You'll want to know: \pa{ਖੰਘ}}
\textbf{\pa{ਖੰਘ}} — \emph{khaṅgh} — \textquotedblleft{}a cough\textquotedblright{}.

\begin{quote}
\pa{ਬੱਚੇ} \pa{ਨੂੰ} \pa{ਰਾਤ} \pa{ਨੂੰ} \pa{ਬਹੁਤ} \pa{ਖੰਘ} \pa{ਸੀ।}
\end{quote}

\begin{quote}
\emph{bachche nū̃ rāt nū̃ bahut khaṅgh sī.}
\end{quote}

\begin{quote}
The child had a bad cough at night.
\end{quote}

A feminine noun.

\begin{cousinweb}
\textbf{\pa{ਖੰਘ}} (\emph{khaṅgh}) goes back to a reconstructed \emph{*khaṅkh}, \textquotedblleft{}cough, hawk\textquotedblright{}, an imitative form; CDIAL lists Punjabi \emph{khaṅgh}, \textquotedblleft{}cough\textquotedblright{}, under it (CDIAL 3763).
\end{cousinweb}

\begin{grammarlens}[title={What you've built}]
One more word for this chapter.
\end{grammarlens}

\subsection*{Your turn}
\begin{itemize}
\raggedright
  \item \emph{Say it:} \emph{khaṅgh}
  \item \emph{Say it:} \emph{khaṅgh}, once more
  \item \emph{Say it:} \emph{bachche nū̃ rāt nū̃ bahut khaṅgh sī.} — The child had a bad cough at night.
  \item \emph{Recall:} say the Punjabi for an injection, then the Punjabi for a head cold, then say \emph{khaṅgh} again
  \item \emph{Write it:} copy \textbf{\pa{ਖੰਘ}} once, checking the whole word
\end{itemize}

\begin{tcolorbox}[breakable,colback=teal!4,colframe=teal!35!black,title={Before you move on}]
What does \pa{ਖੰਘ} mean? (\textquotedblleft{}A cough\textquotedblright{}.) Say it once more.
\end{tcolorbox}
\section[\texorpdfstring{\pa{ਛਿੱਕ} (chhikk)}{ਛਿੱਕ (chhikk)}]{\pa{ਛਿੱਕ} (chhikk) — a sneeze}

\label{lesson:PA-C321-chhikk}



\begin{quote}
Before the new one: say the Punjabi for a head cold, then the Punjabi for a cough.
\end{quote}

\subsection*{You'll want to know: \pa{ਛਿੱਕ}}
\textbf{\pa{ਛਿੱਕ}} — \emph{chhikk} — \textquotedblleft{}a sneeze\textquotedblright{}.

\begin{quote}
\pa{ਮਿਰਚ} \pa{ਤੋਂ} \pa{ਮੈਨੂੰ} \pa{ਛਿੱਕ} \pa{ਆਈ।}
\end{quote}

\begin{quote}
\emph{mirach ton mainū̃ chhikk āī.}
\end{quote}

\begin{quote}
The chilli made me sneeze.
\end{quote}

A feminine noun.

\begin{cousinweb}
\textbf{\pa{ਛਿੱਕ}} (\emph{chhikk}) is inherited from Sanskrit \emph{chikkā}, \textquotedblleft{}a sneeze\textquotedblright{}; CDIAL lists Punjabi \emph{chikk}, \textquotedblleft{}sneeze\textquotedblright{}, under it (CDIAL 5032).
\end{cousinweb}

\begin{grammarlens}[title={What you've built}]
One more word for this chapter.
\end{grammarlens}

\subsection*{Your turn}
\begin{itemize}
\raggedright
  \item \emph{Say it:} \emph{chhikk}
  \item \emph{Say it:} \emph{chhikk}, once more
  \item \emph{Say it:} \emph{mirach ton mainū̃ chhikk āī.} — The chilli made me sneeze.
  \item \emph{Recall:} say the Punjabi for a head cold, then the Punjabi for a cough, then say \emph{chhikk} again
  \item \emph{Write it:} copy \textbf{\pa{ਛਿੱਕ}} once, checking the whole word
\end{itemize}

\begin{tcolorbox}[breakable,colback=teal!4,colframe=teal!35!black,title={Before you move on}]
What does \pa{ਛਿੱਕ} mean? (\textquotedblleft{}A sneeze\textquotedblright{}.) Say it once more.
\end{tcolorbox}
\section[\texorpdfstring{\pa{ਸਾਹ} (sāh)}{ਸਾਹ (sāh)}]{\pa{ਸਾਹ} (sāh) — breath}

\label{lesson:PA-C321-sah}



\begin{quote}
Before the new one: say the Punjabi for a cough, then the Punjabi for a sneeze.
\end{quote}

\subsection*{You'll want to know: \pa{ਸਾਹ}}
\textbf{\pa{ਸਾਹ}} — \emph{sāh} — \textquotedblleft{}breath\textquotedblright{}.

\begin{quote}
\pa{ਪੌੜੀ} \pa{ਤੋਂ} \pa{ਬਾਅਦ} \pa{ਮੇਰਾ} \pa{ਸਾਹ} \pa{ਤੇਜ਼} \pa{ਸੀ।}
\end{quote}

\begin{quote}
\emph{pauṛī ton bāad merā sāh tez sī.}
\end{quote}

\begin{quote}
After the stairs I was breathing fast.
\end{quote}

A masculine noun.

\begin{cousinweb}
\textbf{\pa{ਸਾਹ}} (\emph{sāh}) is inherited from Sanskrit \emph{śvāsa}, \textquotedblleft{}breath, breathing\textquotedblright{}; CDIAL lists Punjabi \emph{sāh}, \textquotedblleft{}breath\textquotedblright{}, under it (CDIAL 12769).
\end{cousinweb}

\begin{grammarlens}[title={What you've built}]
One more word for this chapter.
\end{grammarlens}

\subsection*{Your turn}
\begin{itemize}
\raggedright
  \item \emph{Say it:} \emph{sāh}
  \item \emph{Say it:} \emph{sāh}, once more
  \item \emph{Say it:} \emph{pauṛī ton bāad merā sāh tez sī.} — After the stairs I was breathing fast.
  \item \emph{Recall:} say the Punjabi for a cough, then the Punjabi for a sneeze, then say \emph{sāh} again
  \item \emph{Write it:} copy \textbf{\pa{ਸਾਹ}} once, checking the whole word
\end{itemize}

\begin{tcolorbox}[breakable,colback=teal!4,colframe=teal!35!black,title={Before you move on}]
What does \pa{ਸਾਹ} mean? (\textquotedblleft{}Breath\textquotedblright{}.) Say it once more.
\end{tcolorbox}
\section[\texorpdfstring{\pa{ਨਿਗਲਣਾ} (nigalṇā)}{ਨਿਗਲਣਾ (nigalṇā)}]{\pa{ਨਿਗਲਣਾ} (nigalṇā) — to swallow}

\label{lesson:PA-C321-nigalna}



\begin{quote}
Before the new one: say the Punjabi for a sneeze, then the Punjabi for breath.
\end{quote}

\subsection*{You'll want to know: \pa{ਨਿਗਲਣਾ}}
\textbf{\pa{ਨਿਗਲਣਾ}} — \emph{nigalṇā} — \textquotedblleft{}to swallow\textquotedblright{}.

\begin{quote}
\pa{ਮੈਂ} \pa{ਦਵਾਈ} \pa{ਪਾਣੀ} \pa{ਨਾਲ} \pa{ਨਿਗਲਦੀ} \pa{ਹਾਂ।}
\end{quote}

\begin{quote}
\emph{maiṁ davāī pāṇī nāl nigaldī hāṁ.}
\end{quote}

\begin{quote}
I swallow the medicine with water.
\end{quote}

Past (I …): \pa{ਮੈਂ} \pa{ਦਵਾਈ} \pa{ਨਿਗਲੀ}, \emph{maiṁ davāī niglī}.

\begin{cousinweb}
\textbf{\pa{ਨਿਗਲਣਾ}} (\emph{nigalṇā}) is inherited from Sanskrit \emph{nigalati}, \textquotedblleft{}swallows\textquotedblright{}; CDIAL lists Punjabi \emph{nigalṇā}, \textquotedblleft{}to swallow\textquotedblright{}, under it (CDIAL 7163).
\end{cousinweb}

\begin{grammarlens}[title={What you've built}]
One more word for this chapter.
\end{grammarlens}

\subsection*{Your turn}
\begin{itemize}
\raggedright
  \item \emph{Say it:} \emph{nigalṇā}
  \item \emph{Say it:} \emph{nigalṇā}, once more
  \item \emph{Say it:} \emph{maiṁ davāī pāṇī nāl nigaldī hāṁ.} — I swallow the medicine with water.
  \item \emph{Recall:} say the Punjabi for a sneeze, then the Punjabi for breath, then say \emph{nigalṇā} again
  \item \emph{Write it:} copy \textbf{\pa{ਨਿਗਲਣਾ}} once, checking the whole word
\end{itemize}

\begin{tcolorbox}[breakable,colback=teal!4,colframe=teal!35!black,title={Before you move on}]
What does \pa{ਨਿਗਲਣਾ} mean? (\textquotedblleft{}To swallow\textquotedblright{}.) Say it once more.
\end{tcolorbox}
